\section{One scale sequence for all the estimates}
\label{sec:common-scale-selection}
The estimates are indexed by finite stages, while the final theory
requires arbitrarily high orders. We choose a slowly advancing
stage map so that every fixed stage eventually appears, with enough
time to satisfy its comparison estimates. The same choice also
pays the numerical requests made after the block lengths have
been selected.

\subsection{A slowly advancing sequence with long constant intervals}
Let $b_j,c_j$ be arbitrary nonnegative integers. Define recursively
\[
 A_0=0,\qquad k_j=c_j+2A_j+1,\qquad
                       A_{j+1}=A_j+b_{j+1}+2k_j+1.
\]
The $A_j$ are strictly increasing and tend to infinity. Define
$J_0(m)=j$ when $A_j\le m<A_{j+1}$. These half-open intervals
partition the nonnegative integers. Then $J_0$ is nondecreasing,
$J_0(0)=J_0(1)=0$, and
\[
 J_0(m+1)\le J_0(m)+1,\qquad J_0(m)\longrightarrow\infty.
\]
The increment bound holds because an integer step can cross at
most one of the distinct integer endpoints. The divergence follows
from $m\ge A_j\Rightarrow J_0(m)\ge j$.

If $J_0(m+1)=j>0$, then $A_j\le m+1$ and
$A_j\ge b_j+1$, so $b_j\le m$. Also $k_j\ge c_j$, and the
whole interval
\begin{equation}\label{eq:common-long-plateau}
       k_j-\lfloor k_j/2\rfloor\le m
                        \le k_j+\lfloor k_j/2\rfloor
\end{equation}
lies in $[A_j,A_{j+1})$: its lower endpoint is at least $A_j$
because $k_j\ge2A_j+1$, and its upper endpoint is at most
$2k_j<A_{j+1}$. Thus $J_0=j$ throughout this interval. This
explicit recursion simultaneously postpones entry to every stage
and gives that stage a prescribed long comparison interval.

\subsection{Physical scales, entry requirements and periods}
Let $L_j\ge3$ be the selected nondecreasing odd block lengths,
and let $j_*$ be the first stage. Put
\[
 J(m)=j_*+J_0(m),\qquad s_0=0,\qquad
 s_{m+1}=s_m+\log L_{J(m+1)},\qquad a_m=e^{-s_m}.
\]
Induction and monotonicity give
\begin{equation}\label{eq:common-scale-brackets}
         m\log3\le s_m\le m\log L_{J(m)}.
\end{equation}
In particular $a_m\downarrow0$, and $e^{s_m}$ is the integer
product of the first $m$ block lengths.

In this construction $d_j$ denotes the final entry threshold, including
all the finite-stage enlargements described below. Choose an initial
coupling $\beta_*>1$ with $\beta_*\ge d_{j_*}$, large enough that $t(x)\le1$ for
$x\ge\beta_*$. This is possible since
$t(x)=x^{-1/2}\exp((\log\log x)^2)\to0$. Set
\[
                 \widehat\beta_m=\beta_*+s_Gs_m,
                       \qquad s_G>0.
\]
For each fixed stage $j>j_*$ there are only finitely many entry
and comparison bounds. Require, at sufficiently large $s$,
\begin{equation}\label{eq:common-stage-entries}
 \begin{gathered}
 d_j\le\beta_*+s_Gs,\qquad e_j(\beta_*)\le s,\qquad s>0,\\
 \log L_j\le s/j,\qquad V_j(s)\le e^s.
 \end{gathered}
\end{equation}
Here $d_j$ is the coupling entry, $e_j$ the scale entry, and the
finite-volume requirement is
\[
 V_j(s)=h_j(1+B(s)+B_j^+(s))^{r_j}
             +100(B(s)+1)+100((j+3)L_j+1),
\]
with nonnegative integers $h_j,r_j$ and
\[
 B(s)=\left\lceil\log^2(\beta_*+s_Gs)\right\rceil,
 \qquad
 B_j^+(s)=\left\lceil
       \log^2(\beta_*+s_G(s+\log L_j))\right\rceil.
\]
All inequalities in \eqref{eq:common-stage-entries} hold eventually.
The first four follow from positive linear growth in $s$; for the
last, both ceiling functions are $O(s)$, so $V_j(s)$ is bounded
by a fixed polynomial and hence is $o(e^s)$. Select $b_{j-j_*}$
so that they hold whenever $s\ge b_{j-j_*}\log3$.
The entry property of $J_0$ and the lower bound in
\eqref{eq:common-scale-brackets} now give all these requirements
at every step with $J(m+1)>j_*$. On the first plateau,
$J(m+1)=j_*$ and $\widehat\beta_m\ge\beta_*\ge d_{j_*}$.
Consequently the coupling entrance $d_{J(m+1)}\le\widehat\beta_m$
holds at every step, including $m=0$. The scale-entry and positive-scale
requirements in \eqref{eq:common-stage-entries} are only imposed after
that first plateau; the finite first-plateau volume requirements are paid
by the period choice below.

For the comparison interval at stage $j$, suppose its required
window $w_j(s)$ is nonnegative and integer valued and $w_j(s)/s\to0$, and its
minimum scale is $r_j^0$. Choose $c_{j-j_*}$ so large that
\[
 r_j^0\le s,\qquad w_j(s)\le\lfloor k/2\rfloor
\]
whenever $k\ge c_{j-j_*}$ and
$k\log3\le s\le k\log L_j$. Such a choice exists because
eventually $w_j(s)/s\le(2\log L_j)^{-1}$.
At the center $k=k_{j-j_*}$ of
\eqref{eq:common-long-plateau}, \eqref{eq:common-scale-brackets}
gives precisely these scale bounds. Thus the whole requested
window around $k$ stays at stage $j$.

Only finitely many mesh steps remain in the first plateau.
Choose an integer $M$ at least their finitely many values
$V_{j_*}(s_m)$ and at least the required minimum physical period.
Use periods
\[
                       P_p=4(M+p+1).
\]
The lattice side at level $m$ is $P_p e^{s_m}$. On the first
plateau it pays $V_{j_*}(s_m)$ by the choice of $M$. At every
later step it pays $V_{J(m+1)}(s_m)$ because that quantity is at
most $e^{s_m}$. The three summands of $V_j$ pay respectively the
specified volume, coarse-cover and fine-patch requirements.

The periods increase strictly and have arbitrarily late multiples
of every positive integer $c$. Indeed, given $z$, take
$p=c(M+z+1)-(M+1)$. Then $p\ge z$ and
$P_p=4c(M+z+1)$. Finally
\[
 0\le\frac{\log L_{J(m+1)}}{s_m}\le\frac1{J(m+1)}
                           \longrightarrow0
\]
after the first plateau, by \eqref{eq:common-stage-entries}.
Thus the relative logarithmic mesh increment tends to zero as
required. The stage profiles, block lengths and kernel have all
been retained; the new choices are their scheduling and periods.

\subsection{Countably many requests on that same sequence}
Let $F_{n,i}\in\mathbb R$ be a countable family of numerical
requests selected before constructing the sequence. Define
\[
       E_j=1+\sum_{n=0}^j\sum_{i=0}^j|F_{n,i}|.
\]
Include $\exp(\exp E_j)$ in $d_j$ before choosing $\beta_*$
and applying the preceding construction. Since
$d_{J(m+1)}\le\widehat\beta_m$, monotonicity of the logarithm
twice gives
\[
                 E_{J(m+1)}\le\log\log\widehat\beta_m.
\]
For fixed $n$, the divergence of $J$ supplies a first index after
which $n\le J(m+1)$. At every such $m$ and every $i\le J(m+1)$,
\begin{equation}\label{eq:common-countable-payment}
 |F_{n,i}|\le E_{J(m+1)}
                     \le1+\log\log\widehat\beta_m.
\end{equation}
The first index depends on $n$. We have not required that a finite
stage pay infinitely many independent requests.

There is also an all-stage form. For a prescribed real sequence
$H_j$, include $H_j$ in that same final maximum for $d_j$, again before
choosing $\beta_*$. Then
\begin{equation}\label{eq:common-allstage-payment}
                  H_{J(m+1)}\le\widehat\beta_m
\end{equation}
for every $m$, including zero. Both conclusions persist on every
realization paying enlarged entry thresholds, because their proofs
use only a lower bound on the threshold actually paid. Combining
finitely many such purchases before choosing the realization
therefore retains every earlier conclusion.

\subsection{Selecting the lengths before their numerical costs}
The ordering of choices is essential. First fix the stage profiles,
the original kernel, the row-price families and their secondary
widths. Their analytic suppliers give integer block-size floors.
Take the pointwise maximum of the original floor, the flat and
dual support floors, the curvature floor, the comparison floors,
and the treasury and period floors. Two further families of floors
are named explicitly, because later sections consume them as schedule
conditions.
\begin{enumerate}
\item \emph{Marked-portfolio thresholds.} The marked portfolio,
 its sections and its regradings are fixed before the profiles' price
 maps and the floors
 (Section~\ref{sec:common-observable-purchases}). For that portfolio,
 Section~\ref{sec:ex-marked-portfolio-calibration} produces stage
 thresholds $n^{\mathrm{mk}}_j$ on the block length. Include
 $n^{\mathrm{mk}}_j$ in the floor of stage $j$. Then every block side
 used at stage $j$ is at least $n^{\mathrm{mk}}_j$; this is the
 comparison of thresholds with floor lengths required by the caller
 in Section~\ref{sec:ex-marked-portfolio-calibration}.
\item \emph{Drift and logarithmic-size ratios.} For $j\ge j_*$ let
 $\varepsilon_j\ge0$ be the drift allowance of the stage profile
 $\pi_j$ and $\lambda_j\ge0$ the logarithm of one plus the sum of its
 path allowance, inverse width, column count, clip relaxation and
 support exponent. Both are fixed with the profile. Require
 \begin{equation}\label{eq:common-drift-floor}
  \frac{\varepsilon_j}{\log L_j}\le\min\Bigl\{\frac{s_G}{4},
       \frac1{j^2}\Bigr\},\qquad
  \frac{\lambda_j}{\log L_j}\le\frac1{j^2}\qquad(j\ge j_*).
 \end{equation}
 Since $\varepsilon_j,\lambda_j$ do not depend on the length and
 $c/\log n\to0$ as $n\to\infty$ for each fixed $c\ge0$, there is an
 integer $u_j$ such that \eqref{eq:common-drift-floor} holds for every
 $L_j\ge u_j$; include $u_j$ in the floor of stage $j$.
\end{enumerate}
A nondecreasing odd sequence
$L_j$ above these floors exists: choose each $L_j$ above the
current finite maximum and above $L_{j-1}$. These choices precede
any numerical request that depends on $L_j$.

At step $m$ the block side is $L_{J(m+1)}$ and the incoming stage is
$J(m+1)\ge j_*$. Thus \eqref{eq:common-drift-floor} at $j=J(m+1)$
reads
\[
 \frac{\varepsilon_{J(m+1)}}{\log L_{J(m+1)}}\le
   \min\Bigl\{\frac{s_G}4,\frac1{J(m+1)^2}\Bigr\},
\]
which, together with $J(m)\to\infty$, is the drift condition
\eqref{eq:shoot-drift-ratio} and the second inequality of
\eqref{eq:shoot-adjacent-drift} on the selected schedule, in the
notation of Section~\ref{sec:surviving-interval-shooting} (there
$\varepsilon_{i+1}$ is the allowance of the stage $J(i+1)$ and
$L_i=L_{J(i+1)}$).

With the lengths now fixed, form the actual post-length numerical
costs. If two constructions require $H_j^{(1)}$ and $H_j^{(2)}$,
pay $H_j=\max(H_j^{(1)},H_j^{(2)})$ by
\eqref{eq:common-allstage-payment}. Their countable refinements,
as well as the previously selected inverse-width requests, are
paid by \eqref{eq:common-countable-payment}. Add the nonlinear
entry requirements of Section~\ref{sec:nonlinear-payment-selection}
and every other selected finite-stage threshold before constructing
$J$ and $\widehat\beta$. Every realization paying these enlarged
thresholds preserves all the earlier purchases by the preceding
monotonicity argument. Explicitly, the final entry threshold $d_j$
includes the following stage thresholds, each of which is a finite
maximum at stage $j$ selected before $\beta_*$:
\begin{enumerate}
\item the terminal fee threshold $k_j$ and the entrance delay of the
 marked portfolio, both from
 Section~\ref{sec:common-observable-purchases}; the latter is the delay
 condition of Section~\ref{sec:ex-marked-portfolio-calibration} for the
 adjacent stage pairs, including its profile-cost, credit and tail
 clauses;
\item the startup and reset thresholds and the remaining entry delays
 listed at the end of Section~\ref{sec:common-funded-history},
 selected after the protection parameters there;
\item the loop-capacity threshold $\Delta_j$ of
 Lemma~\ref{lem:common-loop-capacity} below, for the side-dependent
 thresholds $W_j(L_j)$ of Section~\ref{sec:path-attachment-producer}.
\end{enumerate}
Since $d_{J(m+1)}\le\widehat\beta_m$ at every step, every one of these
thresholds is paid at stage $J(m+1)$ by the coupling $\widehat\beta_m$.
In particular $k_{J(m+1)}\le\widehat\beta_m$, which is the fee condition
used in Proposition~\ref{prop:terminal-marked-voice}.

\begin{lemma}[Loop capacities]\label{lem:common-loop-capacity}
Let $\theta_j>0$ be nonincreasing protection exponents, fixed
after the block lengths, and put
\[
 P(x)=\exp\bigl((\log\log x)^2\bigr),\qquad
 \mathrm{loopCap}_m=2P(\widehat\beta_m)^{\theta_{J(m)}}.
\]
For every real sequence $W_j$ there are thresholds $\Delta_j$ such that,
whenever $\Delta_j\le d_j$ for all $j$, the selected schedule satisfies
at every $m\ge0$, with $j=J(m+1)$,
\begin{equation}\label{eq:common-loop-capacity}
 W_j\le\mathrm{loopCap}_m,\qquad
 \tfrac12\mathrm{loopCap}_{m+1}\le\mathrm{loopCap}_m,\qquad
 1\le\mathrm{loopCap}_m.
\end{equation}
\end{lemma}
\begin{proof}
Fix $j$ and put $v_j=s_G\log L_j\ge0$. As $x\to\infty$,
$2P(x)^{\theta_j}\to\infty$, so eventually $W_j\le2P(x)^{\theta_j}$.
For the second requirement let $\ell=\log x$ and $\ell'=\log(x+v_j)$.
For $x>1$, $\ell\le\ell'\le\ell+v_j/x$, since
$\log(1+v_j/x)\le v_j/x$. Using $\log t\le t-1$,
\[
 0\le\log\ell'-\log\ell\le\frac{\ell'}{\ell}-1\le\frac{v_j}{x\ell}.
\]
Once $\ell\ge1$ and $x\ge v_j$, also $0\le\log\ell\le\ell$ and
$\log\ell'\le\ell'-1\le\ell+v_j/x-1\le\ell+1$, so
\[
 (\log\ell')^2-(\log\ell)^2
   \le(\log\ell'+\log\ell)\frac{v_j}{x\ell}
   \le\frac{(2\ell+1)v_j}{x\ell}\le\frac{3v_j}{x}.
\]
Hence $P(x+v_j)^{\theta_j}\le e^{3v_j\theta_j/x}P(x)^{\theta_j}
\le2P(x)^{\theta_j}$ as soon as $3v_j\theta_j/x\le\log2$. Let
$\Delta_j$ be a common threshold beyond which both properties hold,
and $x\ge\max\{e,v_j\}$.

Now let $m\ge0$, $j=J(m+1)$ and $x=\widehat\beta_m\ge d_j\ge\Delta_j$.
Since $P\ge1$ and $\theta_j\le\theta_{J(m)}$ (because $J(m)\le j$ and
the exponents are nonincreasing),
$P(x)^{\theta_j}\le P(x)^{\theta_{J(m)}}$. The first property gives
$W_j\le\mathrm{loopCap}_m$. Next
$\widehat\beta_{m+1}=x+s_G\log L_{J(m+1)}=x+v_j$ and $J(m+1)=j$, so
\[
 \tfrac12\mathrm{loopCap}_{m+1}=P(x+v_j)^{\theta_j}
   \le2P(x)^{\theta_j}\le2P(x)^{\theta_{J(m)}}=\mathrm{loopCap}_m.
\]
Finally $P(x)\ge1$ and $\theta_{J(m)}>0$ give $\mathrm{loopCap}_m\ge2$.
\end{proof}
These are the three loop-capacity inequalities used in
Section~\ref{sec:path-attachment-producer}, with $W_j$ its side-dependent
threshold at stage $j$ evaluated at $L_j$; the threshold is indexed by
the incoming stage and the capacity by the output level, as there.

Every step then uses the original length $L_{J(m+1)}$, so its
actual finite patch is exactly the patch for which the post-length
cost was computed. In particular a bound of the form
$e^{H_j}\le\widehat\beta_m$ yields
$H_j\le\log\widehat\beta_m$ on that patch; no estimate for a
different block can be substituted. A maximum floor gives each
of its component floors, and a maximum post-length cost gives
each component cost. Applying the native analytic suppliers to
these inequalities therefore supplies their conclusions on the
same original levels and quadratic horizon.

This proves the scheduling and numerical selection step. The
analytic hypotheses of those suppliers, and the sublinear window
bounds used above, are still the specific estimates established
in their own sections. Numerical selection does not replace them.
